\ifdefined\BookMaster
\def\BookEnd{}
\else
\documentclass[a4paper,12pt,fleqn,openany,twoside,fontset=windows]{ctexbook}

\usepackage[fleqn]{amsmath}
\usepackage{amssymb,amsthm,mathrsfs}
\usepackage{geometry}
\usepackage[dvipsnames]{xcolor}
\usepackage{graphicx}
\usepackage{tikz}
\usetikzlibrary{calc}
\usepackage[most]{tcolorbox}
\usepackage{fancyhdr}
\usepackage{titlesec}
\usepackage{tocloft}
\usepackage{enumitem}
\usepackage{microtype}
\usepackage{pifont}
\usepackage{hyperref}
\usepackage{romannum}
\usepackage{mathtools}

\definecolor{bookblue}{HTML}{264653}
\definecolor{bookteal}{HTML}{4F7C82}
\definecolor{booklight}{HTML}{EDF4F3}
\definecolor{bookgray}{HTML}{5E686B}

\geometry{
    inner=2.7cm,
    outer=2.5cm,
    top=2.7cm,
    bottom=2.7cm,
    headheight=18pt,
    headsep=18pt,
    footskip=25pt
}
\setlength{\mathindent}{0pt}
\setlength{\parindent}{5pt}
\setlength{\parskip}{3pt plus 1pt minus 1pt}
\linespread{1.25}
\allowdisplaybreaks[3]
\AtBeginDocument{%
    \setlength{\parindent}{5pt}
    \setlength{\abovedisplayskip}{8pt plus 2pt minus 2pt}
    \setlength{\belowdisplayskip}{8pt plus 2pt minus 2pt}
    \setlength{\abovedisplayshortskip}{5pt plus 1pt minus 1pt}
    \setlength{\belowdisplayshortskip}{5pt plus 1pt minus 1pt}
}

\hypersetup{
    unicode=true,
    colorlinks=true,
    linkcolor=bookblue,
    citecolor=bookblue,
    urlcolor=bookteal,
    bookmarksopen=true,
    pdfauthor={Chen},
    pdftitle={数学分析习题整理}
}

\renewcommand{\chaptermark}[1]{%
    \edef\@tmp{第\thechapter 章\quad #1}%
    \expandafter\markboth\expandafter{\@tmp}{}%
}
\fancypagestyle{main}{%
    \fancyhf{}
    \fancyhead[LE]{\small\color{bookgray}\nouppercase{\leftmark}}
    \fancyhead[RO]{\small\color{bookgray}数学分析习题整理}
    \fancyfoot[C]{\color{bookblue}\thepage}
    \fancyfoot[R]{\small\bfseries Chen\;\; github.com/mathlover-1/math-analysis}
    \renewcommand{\headrulewidth}{0.4pt}
    \renewcommand{\footrulewidth}{0pt}
}
\fancypagestyle{plain}{%
    \fancyhf{}
    \fancyfoot[C]{\color{bookblue}\thepage}
    \fancyfoot[R]{\small\bfseries Chen\;\; github.com/mathlover-1/math-analysis}
    \renewcommand{\headrulewidth}{0pt}
    \renewcommand{\footrulewidth}{0pt}
}
\pagestyle{main}

\titleformat{\chapter}[display]
  {\normalfont\sffamily\bfseries\color{bookblue}}
  {\raggedleft\fontsize{46}{50}\selectfont\thechapter}
  {-18pt}
  {\titlerule[1.2pt]\vspace{10pt}\filright\Huge}
\titlespacing*{\chapter}{0pt}{-15pt}{36pt}

\renewcommand{\contentsname}{目\quad 录}
\renewcommand{\cfttoctitlefont}{\hfill\sffamily\bfseries\Huge\color{bookblue}}
\renewcommand{\cftaftertoctitle}{\hfill}
\setlength{\cftbeforetoctitleskip}{5pt}
\setlength{\cftaftertoctitleskip}{28pt}
\setlength{\cftbeforechapskip}{12pt}
\renewcommand{\cftchapfont}{\sffamily\bfseries\large\color{bookblue}}
\renewcommand{\cftchappagefont}{\sffamily\bfseries\large\color{bookblue}}
\renewcommand{\cftchapleader}{\color{bookteal}\cftdotfill{2.5}}

\newtcolorbox{ti}{
    enhanced,
    breakable,
    colback=booklight!55!white,
    colframe=bookteal!35!white,
    boxrule=0.45pt,
    boxsep=0pt,
    left=10pt,
    right=10pt,
    top=9pt,
    bottom=9pt,
    before skip=14pt,
    after skip=14pt,
    arc=2pt,
    outer arc=2pt,
    before upper={\parindent=0pt}
}

\newenvironment{booknotebody}{%
    \begin{center}
    \begin{minipage}{0.84\textwidth}
    \songti\zihao{-4}\linespread{1.45}\selectfont
    \setlength{\parindent}{2\ccwd}
    \setlength{\parskip}{0.75em}
}{%
    \end{minipage}
    \end{center}
}

\newenvironment{booknotelongbody}{%
    \begingroup
    \songti\zihao{-4}\linespread{1.45}\selectfont
    \setlength{\parindent}{2\ccwd}
    \setlength{\parskip}{0.75em}
    \begin{list}{}{%
        \setlength{\leftmargin}{0.08\textwidth}
        \setlength{\rightmargin}{0.08\textwidth}
        \setlength{\topsep}{0pt}
        \setlength{\partopsep}{0pt}
        \setlength{\parsep}{0pt}
        \setlength{\itemsep}{0pt}
    }
    \item\relax
}{%
    \end{list}
    \endgroup
}

\fancypagestyle{plainnowm}{%
    \fancyhf{}
    \fancyfoot[C]{\color{bookblue}\thepage}
    \renewcommand{\headrulewidth}{0pt}
    \renewcommand{\footrulewidth}{0pt}
}

\newcommand{\booknotetitle}[2][plain]{%
    \clearpage
    \phantomsection
    \addcontentsline{toc}{chapter}{#2}
    \markboth{#2}{}
    \thispagestyle{#1}
    \vspace*{1.2cm}
    \begin{center}
        {\sffamily\heiti\bfseries\fontsize{26}{34}\selectfont\color{bookblue}#2\par}
        \vspace{0.8em}
        {\color{bookteal}\rule{4.5cm}{0.8pt}\par}
    \end{center}
    \vspace{0.6cm}
}

\renewcommand{\proofname}{证明}
\renewcommand{\qedsymbol}{\color{bookteal}\ensuremath{\blacksquare}}

\title{数学分析习题整理}
\author{Chen}
\date{}

\makeatletter
% ==================== 旧版封面/封底设计备份 ====================
% 2026-08 重设计,旧代码用 \iffalse...\fi 封存,新版见下方。
\iffalse
\renewcommand{\maketitle}{%
    \begin{titlepage}
        \thispagestyle{empty}
        \setcounter{page}{0}
        \noindent\hspace*{-\parindent}\begin{tikzpicture}[remember picture,overlay]
            \fill[bookblue] (current page.north west) rectangle ([yshift=-4.2cm]current page.north east);
            \fill[booklight] ([yshift=-4.2cm]current page.north west) rectangle (current page.south east);
            \draw[bookteal!55,line width=0.7pt]
                ([xshift=2.4cm,yshift=4.3cm]current page.south west)
                .. controls +(2.5cm,1.8cm) and +(-2.8cm,-1.3cm) ..
                ([xshift=-2.4cm,yshift=7.2cm]current page.south east);
            \draw[bookteal!35,line width=0.7pt]
                ([xshift=2.5cm,yshift=6.2cm]current page.south west)
                .. controls +(3.0cm,-1.2cm) and +(-3.0cm,1.6cm) ..
                ([xshift=-2.3cm,yshift=3.8cm]current page.south east);
        \end{tikzpicture}
        \vspace*{5.6cm}
        \noindent\color{bookblue}\rule{\textwidth}{1.4pt}\par
        \vspace{8pt}
        {\centering\sffamily\bfseries\fontsize{32}{42}\selectfont\color{bookblue}\@title\par}
        \vspace{8pt}
        \noindent\color{bookteal}\rule{\textwidth}{0.6pt}\par
        \vspace{2.6cm}
        {\centering\Large\color{bookgray}\@author\par}
        \vfill
        {\centering\color{bookteal}\large
            $\displaystyle \lim_{n\to\infty} a_n$\qquad
            $\displaystyle \int_a^b f(x)\,\mathrm{d}x$\qquad
            $\displaystyle \nabla f(x,y)$\par
        }
        \vspace*{1.5cm}
    \end{titlepage}
}

\newcommand{\makebackcover}{%
    \begingroup
    \let\cleardoublepage\clearpage
    \begin{titlepage}
        \thispagestyle{empty}
        \noindent\hspace*{-\parindent}\begin{tikzpicture}[remember picture,overlay]
            \fill[bookblue] (current page.north west) rectangle (current page.south east);
            \fill[booklight,rounded corners=7pt]
                ([xshift=2.2cm,yshift=-3.0cm]current page.north west)
                rectangle
                ([xshift=-2.2cm,yshift=3.0cm]current page.south east);
            \fill[bookteal!80]
                (current page.north west) --
                ([xshift=5.3cm]current page.north west) --
                ([xshift=2.1cm,yshift=-4.4cm]current page.north west) --
                ([yshift=-6.2cm]current page.north west) -- cycle;
            \fill[bookteal!55]
                (current page.south east) --
                ([xshift=-5.8cm]current page.south east) --
                ([xshift=-2.4cm,yshift=4.2cm]current page.south east) --
                ([yshift=6.0cm]current page.south east) -- cycle;
            \foreach \radius in {0.75,1.15,1.55,1.95} {
                \draw[bookteal!55,line width=0.55pt]
                    ([xshift=-2.9cm,yshift=-4.4cm]current page.north east)
                    circle[radius=\radius cm];
            }
            \foreach \radius in {0.65,1.05,1.45} {
                \draw[booklight!65,line width=0.55pt]
                    ([xshift=2.4cm,yshift=3.6cm]current page.south west)
                    circle[radius=\radius cm];
            }
            \fill[bookteal!16,rounded corners=3pt]
                ([xshift=3.5cm,yshift=-7.1cm]current page.north west)
                rectangle
                ([xshift=-5.0cm,yshift=7.2cm]current page.south east);
            \fill[bookteal!30,rounded corners=3pt]
                ([xshift=4.4cm,yshift=-8.0cm]current page.north west)
                rectangle
                ([xshift=-4.1cm,yshift=8.1cm]current page.south east);
            \begin{scope}[shift={([xshift=-0.75cm]current page.center)}]
                \clip[rounded corners=3pt] (-5.25,-4.75) rectangle (5.25,4.75);
                \foreach \v in {-3.6,-3.0,...,3.6} {
                    \draw[bookteal!48,line width=0.45pt,opacity=0.72]
                        plot[domain=-3.6:3.6,samples=45,smooth,variable=\u]
                        ({\u+0.55*\v},{0.28*\u-0.35*\v+0.11*(\u*\u-\v*\v)});
                }
                \foreach \u in {-3.6,-3.0,...,3.6} {
                    \draw[bookblue!42,line width=0.42pt,opacity=0.62]
                        plot[domain=-3.6:3.6,samples=45,smooth,variable=\v]
                        ({\u+0.55*\v},{0.28*\u-0.35*\v+0.11*(\u*\u-\v*\v)});
                }
                \draw[bookteal!70,line width=0.75pt,opacity=0.82]
                    plot[domain=-3.6:3.6,samples=55,smooth,variable=\u]
                    ({\u},{0.28*\u+0.11*\u*\u});
                \draw[bookblue!60,line width=0.7pt,opacity=0.75]
                    plot[domain=-3.6:3.6,samples=55,smooth,variable=\v]
                    ({0.55*\v},{-0.35*\v-0.11*\v*\v});
            \end{scope}
            \draw[bookteal!45,line width=0.65pt]
                ([xshift=3.0cm,yshift=-8.2cm]current page.north west)
                .. controls +(2.2cm,2.0cm) and +(-2.0cm,-2.1cm) ..
                ([xshift=-3.0cm,yshift=9.0cm]current page.south east);
            \draw[bookteal!25,line width=0.55pt]
                ([xshift=3.1cm,yshift=8.0cm]current page.south west)
                .. controls +(2.7cm,-1.7cm) and +(-2.6cm,1.8cm) ..
                ([xshift=-3.1cm,yshift=-8.7cm]current page.north east);
            \node[anchor=west,text=bookteal!72,scale=0.78] at
                ([xshift=2.8cm,yshift=-4.0cm]current page.north west) {%
                $\displaystyle e^{\mathrm{i}\pi}+1=0$};
            \node[anchor=east,text=bookblue!62,scale=0.67] at
                ([xshift=-2.7cm,yshift=4.0cm]current page.south east) {%
                $\displaystyle
                f(x)=\sum_{n=0}^{\infty}\frac{f^{(n)}(a)}{n!}(x-a)^n$};
        \end{tikzpicture}
        \mbox{}
    \end{titlepage}
    \endgroup
}
\fi

% ==================== 新版封面(2026-08 重设计) ====================
% 设计要点:顶部深绿纵向渐变+波浪下缘+两层半透明过渡波(解决深浅绿生硬过渡);
% 深色区螺旋线/同心圆母题;标题+英文副题+作者居中;
% 中部 Fourier 部分和逼近方波曲线;底部浅色波纹(无文字)。
\renewcommand{\maketitle}{%
    \begin{titlepage}
        \thispagestyle{empty}
        \setcounter{page}{0}
        \noindent\hspace*{-\parindent}\begin{tikzpicture}[remember picture,overlay]
            % ---------- 底色 ----------
            \fill[booklight] (current page.north west) rectangle (current page.south east);

            % ---------- 顶部深色区:纵向渐变 + 波浪下缘 ----------
            \shade[top color=bookblue, bottom color=bookteal!55!bookblue]
                (current page.north west) -- (current page.north east) --
                ([yshift=-8.6cm]current page.north east)
                .. controls ([xshift=-5.2cm,yshift=-10.6cm]current page.north east)
                         and ([xshift=6.0cm,yshift=-8.0cm]current page.north west) ..
                ([yshift=-10.0cm]current page.north west) -- cycle;

            % ---------- 过渡层 1:半透明波浪 ----------
            \fill[bookteal!60!bookblue,opacity=0.42]
                ([yshift=-8.6cm]current page.north east)
                .. controls ([xshift=-5.2cm,yshift=-10.6cm]current page.north east)
                         and ([xshift=6.0cm,yshift=-8.0cm]current page.north west) ..
                ([yshift=-10.0cm]current page.north west) --
                ([yshift=-12.6cm]current page.north west)
                .. controls ([xshift=6.5cm,yshift=-11.0cm]current page.north west)
                         and ([xshift=-5.5cm,yshift=-13.2cm]current page.north east) ..
                ([yshift=-11.4cm]current page.north east) -- cycle;

            % ---------- 过渡层 2:更浅的波浪 ----------
            \fill[bookteal!35,opacity=0.35]
                ([yshift=-10.6cm]current page.north east)
                .. controls ([xshift=-6.0cm,yshift=-12.4cm]current page.north east)
                         and ([xshift=5.0cm,yshift=-9.8cm]current page.north west) ..
                ([yshift=-11.8cm]current page.north west) --
                ([yshift=-13.8cm]current page.north west)
                .. controls ([xshift=5.5cm,yshift=-12.6cm]current page.north west)
                         and ([xshift=-6.5cm,yshift=-14.6cm]current page.north east) ..
                ([yshift=-12.8cm]current page.north east) -- cycle;

            % ---------- 深色区装饰:右上螺旋线 ----------
            \begin{scope}[shift={([xshift=-3.4cm,yshift=-4.3cm]current page.north east)}]
                \draw[booklight!45!bookteal,opacity=0.55,line width=0.6pt]
                    plot[domain=0:1440,samples=220,smooth,variable=\t]
                    ({0.0028*\t*cos(\t)},{0.0028*\t*sin(\t)});
            \end{scope}

            % ---------- 深色区装饰:左下同心圆 ----------
            \foreach \r in {0.7,1.25,1.8,2.35,2.9}{
                \draw[booklight!35!bookteal,opacity=0.5,line width=0.5pt]
                    ([xshift=0.4cm,yshift=-7.9cm]current page.north west) circle[radius=\r cm];
            }

            % ---------- 主标题 ----------
            \node[anchor=center,text=booklight] at
                ([yshift=-4.5cm]current page.north)
                {\sffamily\bfseries\fontsize{40}{46}\selectfont \@title};

            % ---------- 分隔短线 ----------
            \draw[booklight!65,line width=1.2pt]
                ([xshift=-2.6cm,yshift=-5.75cm]current page.north) --
                ([xshift=2.6cm,yshift=-5.75cm]current page.north);

            % ---------- 英文副题 ----------
            \node[anchor=center,text=booklight!70!bookteal] at
                ([yshift=-6.55cm]current page.north)
                {\itshape\large Mathematical Analysis \textperiodcentered\
                 A Collection of Exercises};

            % ---------- 作者(紧随标题下方) ----------
            \node[anchor=center,text=booklight!88] at
                ([yshift=-8.2cm]current page.north)
                {\sffamily\large \@author\quad 整理};

            % ---------- 中部:Fourier 部分和逼近方波 ----------
            \begin{scope}[shift={([xshift=-5.2cm,yshift=-18.1cm]current page.north)}]
                % 坐标轴
                \draw[bookgray!55,line width=0.5pt,->] (-0.5,0) -- (10.9,0);
                \draw[bookgray!55,line width=0.5pt,->] (0,-1.75) -- (0,1.95);
                % 目标方波(虚线)
                \draw[bookgray!75,line width=0.6pt,dashed]
                    (0,0.92) -- (5,0.92) (5,-0.92) -- (10,-0.92);
                % N=1
                \draw[bookteal!45,line width=0.7pt]
                    plot[domain=0:360,samples=90,smooth,variable=\x]
                    ({\x/36},{1.18*sin(\x)});
                % N=3
                \draw[bookteal!80,line width=0.8pt]
                    plot[domain=0:360,samples=120,smooth,variable=\x]
                    ({\x/36},{1.18*(sin(\x)+sin(3*\x)/3)});
                % N=7
                \draw[bookblue,line width=0.95pt]
                    plot[domain=0:360,samples=150,smooth,variable=\x]
                    ({\x/36},{1.18*(sin(\x)+sin(3*\x)/3+sin(5*\x)/5+sin(7*\x)/7)});
                % 刻度标注
                \node[anchor=north east,text=bookgray!80,scale=0.62] at (-0.12,0) {$0$};
                \node[anchor=north,text=bookgray!80,scale=0.62] at (5,0) {$\pi$};
                \node[anchor=north,text=bookgray!80,scale=0.62] at (10,0) {$2\pi$};
            \end{scope}
            % 图注
            \node[anchor=center,text=bookgray,scale=0.8] at
                ([yshift=-20.5cm]current page.north)
                {$S_N(x)=\displaystyle\sum_{k=1}^{N}\frac{\sin\bigl((2k-1)x\bigr)}{2k-1}$
                 \ 逐步逼近方波};

            % ---------- 底部浅色波纹(无文字,与顶部波浪呼应) ----------
            \fill[bookteal!22,opacity=0.5]
                ([yshift=2.4cm]current page.south west)
                .. controls ([xshift=6.0cm,yshift=3.9cm]current page.south west)
                         and ([xshift=-6.5cm,yshift=1.0cm]current page.south east) ..
                ([yshift=2.9cm]current page.south east) --
                (current page.south east) -- (current page.south west) -- cycle;
            \fill[bookteal!40!bookblue,opacity=0.28]
                ([yshift=1.3cm]current page.south west)
                .. controls ([xshift=5.5cm,yshift=2.7cm]current page.south west)
                         and ([xshift=-5.5cm,yshift=0.3cm]current page.south east) ..
                ([yshift=1.9cm]current page.south east) --
                (current page.south east) -- (current page.south west) -- cycle;
        \end{tikzpicture}
        \mbox{}
    \end{titlepage}
}

% ==================== 新版封底(纯背景,2026-08 重设计) ====================
% 设计要点:对角渐变底+角部半透明叠层;中下部外摆线(玫瑰形)为视觉主体;
% 底部两层半透明白色波浪;无任何文字。
\newcommand{\makebackcover}{%
    \begingroup
    \let\cleardoublepage\clearpage
    \begin{titlepage}
        \thispagestyle{empty}
        \noindent\hspace*{-\parindent}\begin{tikzpicture}[remember picture,overlay]
            % ---------- 底色:对角渐变 ----------
            \shade[top color=bookblue, bottom color=bookteal!48!bookblue,
                   shading angle=-30]
                (current page.north west) rectangle (current page.south east);

            % ---------- 角部半透明叠层 ----------
            \fill[white,opacity=0.05]
                ([xshift=2.0cm,yshift=1.0cm]current page.north east) circle[radius=7.2cm];
            \fill[white,opacity=0.06]
                ([xshift=-0.5cm,yshift=-1.5cm]current page.north east) circle[radius=3.6cm];

            % ---------- 左下同心圆 ----------
            \foreach \r in {1.0,1.75,2.5,3.25,4.0,4.75}{
                \draw[booklight!40,opacity=0.55,line width=0.55pt]
                    ([xshift=-0.6cm,yshift=-0.4cm]current page.south west) circle[radius=\r cm];
            }
            % ---------- 右下螺旋线 ----------
            \begin{scope}[shift={([xshift=-2.9cm,yshift=3.2cm]current page.south east)}]
                \draw[booklight!45!bookteal,opacity=0.5,line width=0.55pt]
                    plot[domain=0:1260,samples=200,smooth,variable=\t]
                    ({0.0036*\t*cos(\t)},{0.0036*\t*sin(\t)});
            \end{scope}
            % ---------- 左上斜线组 ----------
            \foreach \i in {0,1,2}{
                \draw[bookteal!60!booklight,opacity=0.5,line width=0.6pt]
                    ([xshift={1.6cm+\i*0.45cm}]current page.north west) --
                    ([yshift={-1.6cm-\i*0.45cm}]current page.north west);
            }

            % ---------- 中下部视觉主体:外摆线(玫瑰形) ----------
            \begin{scope}[shift={([yshift=-16.6cm]current page.north)}]
                % 外圈
                \draw[booklight!35,opacity=0.6,line width=0.6pt]
                    (0,0) circle[radius=4.9cm];
                \draw[booklight!22,opacity=0.5,line width=0.5pt]
                    (0,0) circle[radius=5.25cm];
                % 主玫瑰线
                \draw[booklight!55!bookteal,opacity=0.8,line width=0.85pt]
                    plot[domain=0:1080,samples=400,smooth,variable=\t]
                    ({0.335*(8*cos(\t)-5*cos(2.6667*\t))},
                     {0.335*(8*sin(\t)-5*sin(2.6667*\t))});
                % 内层玫瑰线
                \draw[booklight!38!bookteal,opacity=0.6,line width=0.6pt]
                    plot[domain=0:1080,samples=400,smooth,variable=\t]
                    ({0.21*(8*cos(\t)-5*cos(2.6667*\t))},
                     {0.21*(8*sin(\t)-5*sin(2.6667*\t))});
            \end{scope}

            % ---------- 底部半透明波浪 ----------
            \fill[white,opacity=0.06]
                ([yshift=4.6cm]current page.south west)
                .. controls ([xshift=6.5cm,yshift=6.4cm]current page.south west)
                         and ([xshift=-6.0cm,yshift=2.6cm]current page.south east) ..
                ([yshift=5.2cm]current page.south east) --
                (current page.south east) -- (current page.south west) -- cycle;
            \fill[booklight,opacity=0.08]
                ([yshift=2.6cm]current page.south west)
                .. controls ([xshift=5.5cm,yshift=4.3cm]current page.south west)
                         and ([xshift=-6.5cm,yshift=1.0cm]current page.south east) ..
                ([yshift=3.3cm]current page.south east) --
                (current page.south east) -- (current page.south west) -- cycle;
        \end{tikzpicture}
        \mbox{}
    \end{titlepage}
    \endgroup
}
\makeatother

% 页脚右下方常态化水印(署名 + 仓库地址,加粗黑色),已在 main/plain 分页样式中定义,所有页面均生效.

\begin{document}
\mainmatter
\pagestyle{main}
\setcounter{chapter}{7}
\def\BookEnd{\end{document}}
\fi

% 子标题编号格式:§1,§2,§3,...(带段落符号);标题居左.
\renewcommand{\thesection}{\S\arabic{section}}
\setcounter{section}{0}
\titleformat{\section}[block]
  {\normalfont\heiti\bfseries\Large\color{bookblue}}
  {\thesection}{1em}{}
\titlespacing*{\section}{0pt}{3.2em plus 0.4em}{1em}

% 子节(subsection)编号格式:§X.Y(父节号.子节号),标题居左,比小节略小.
\renewcommand{\thesubsection}{\S\arabic{section}.\arabic{subsection}}
\titleformat{\subsection}[block]
  {\normalfont\heiti\bfseries\large\color{bookblue}}
  {\thesubsection}{1em}{}
\titlespacing*{\subsection}{0pt}{2.4em plus 0.3em}{0.8em}

% 附录行间编号公式居中:正文已按 fleqn(左对齐)排版,本附录要求 equation/align 等编号公式居中,
% 故在此关闭 amsmath 的 fleqn 布局分支(无编号 $$...$$ 公式本就居中,不受影响)。仅作用于本附录及其后内容。
\makeatletter\@fleqnfalse\makeatother

% 附录公式编号:去掉"章.节"前缀,同一节内从 1 开始递推(1,2,3,...)。
\makeatletter
\@addtoreset{equation}{section}
\renewcommand{\theequation}{\arabic{equation}}
\makeatother

% 附录内局部放大绿框(ti)之间的间距:+20%(14pt → 16.8pt),不影响正文各章.
\renewtcolorbox{ti}{
    enhanced,
    breakable,
    colback=booklight!55!white,
    colframe=bookteal!35!white,
    boxrule=0.45pt,
    boxsep=0pt,
    left=10pt,
    right=10pt,
    top=9pt,
    bottom=9pt,
    before skip=16.8pt,
    after skip=16.8pt,
    arc=2pt,
    outer arc=2pt,
    before upper={\parindent=0pt}
}
\setcounter{section}{6}   % 从§7开始编号
% ==================== 附录:数学分析重要定理(第七至八节) ====================
% 本文件为附录的后段:§7 级数 与 §8 Fourier分析.
% 前两段分别见 附录-数学分析重要定理-第1至4节.tex 和 附录-数学分析重要定理-第5至6节.tex.
% 三文件均可独立编译;主文件中按顺序 input 三者即可.
% ========================================================================
\section{级数}
\subsection{数项级数}
\begin{ti}
\textbf{定义7.1(级数的收敛与发散)}  设$\{a_n\}$为数列,称形式和$\displaystyle\sum_{n=1}^{\infty}a_n=a_1+a_2+\cdots+a_n+\cdots$为\textbf{无穷级数},
$a_n$称为通项或一般项,$S_n=\displaystyle\sum_{k=1}^{n}a_k$称为级数的第$n$个\textbf{部分和}.
如果数列$\{S_n\}$收敛,则称级数$\displaystyle\sum_{n=1}^{\infty}a_n$\textbf{收敛},其和记为$\displaystyle\sum_{n=1}^{\infty}a_n=\lim_{n\to\infty}S_n$;否则就称级数$\displaystyle\sum_{n=1}^{\infty}a_n$\textbf{发散}.
\end{ti}

\begin{ti}
\textbf{命题7.2(收敛的必要条件)}  如果级数$\displaystyle\sum_{n=1}^{\infty}a_n$收敛,则$\displaystyle\lim_{n\to\infty}a_n=0$.
\end{ti}

\begin{ti}
\textbf{引理7.3(级数与广义积分)}  设$\{a_n\}$为一列实数,在$[1,+\infty)$中定义函数$a(x)$如下:设$k\geq 1$,当$x\in[k,k+1)$时,令$a(x)=a_k$.则级数$\displaystyle\sum_{n=1}^{\infty}a_n$收敛当且仅当无穷积分$\displaystyle\int_1^{+\infty}a(x)\,\mathrm{d}x$收敛,且收敛时级数的和等于无穷积分的值.
\end{ti}

\begin{ti}
\textbf{定理7.4($Dirichlet$)}  设数列$\{a_n\}$单调地趋于$0$,级数$\displaystyle\sum_{n=1}^{\infty}b_n$的部分和有界,则级数$\displaystyle\sum_{n=1}^{\infty}a_nb_n$收敛.
\end{ti}

\begin{ti}
\textbf{定理7.5($Abel$)}  设$\{a_n\}$为有界的单调数列,$\displaystyle\sum_{n=1}^{\infty}b_n$收敛,则级数$\displaystyle\sum_{n=1}^{\infty}a_nb_n$收敛.
\end{ti}

\begin{ti}
\textbf{推论7.6($Leibniz$)}  设$a_n$单调地趋于$0$,则级数$\displaystyle\sum_{n=1}^{\infty}(-1)^{n}a_n$收敛.
\end{ti}

\begin{ti}
\textbf{命题7.7($Abel$变换)}  设$\{a_i\}$,$\{b_i\}$为两组实数,如果$b_i=B_i-B_{i-1}$,则
$$\sum_{i=1}^{n}a_i b_i=\sum_{i=1}^{n-1}(a_i-a_{i+1})B_i+a_n B_n-a_1 B_0.$$
\end{ti}

\begin{ti}
\textbf{定义7.8(绝对收敛与条件收敛)}  如果级数$\displaystyle\sum_{n=1}^{\infty}\left\lvert a_n\right\rvert$收敛,则称$\displaystyle\sum_{n=1}^{\infty}a_n$\textbf{绝对收敛};
如果$\displaystyle\sum_{n=1}^{\infty}a_n$收敛而$\displaystyle\sum_{n=1}^{\infty}\left\lvert a_n\right\rvert$发散,则称其\textbf{条件收敛}.
\end{ti}

\begin{ti}
\textbf{定理7.9(比较判别法)}  设$\displaystyle\sum a_n$和$\displaystyle\sum b_n$为正项级数,如果存在$M>0,N\geq 1$,使得当$n>N$时$a_n\leq Mb_n$,则$\displaystyle\sum b_n$收敛时$\displaystyle\sum a_n$也收敛;$\displaystyle\sum a_n$发散时$\displaystyle\sum b_n$也发散.
\end{ti}

\begin{ti}
\textbf{定理7.10(比值判别法)}  设$a_n>0$,且$\displaystyle\lim_{n\to\infty}\dfrac{a_{n+1}}{a_n}=r$.\\
(1)当$r<1$时级数$\displaystyle\sum a_n$收敛;\\
(2)当$r>1$时级数$\displaystyle\sum a_n$发散.
\end{ti}

\begin{ti}
\textbf{定理7.11(根值判别法)}  设$a_n\geq 0$,且$\displaystyle\lim_{n\to\infty}\sqrt[n]{a_n}=r$.\\
(1)当$r<1$时级数$\displaystyle\sum a_n$收敛;\\
(2)当$r>1$时级数$\displaystyle\sum a_n$发散.
\end{ti}

\begin{ti}
\textbf{定理7.12(积分判别法)}  设函数$f$在$[1,+\infty)$上非负单调递减,则级数$\displaystyle\sum_{n=1}^{\infty}f(n)$与广义积分$\displaystyle\int_1^{+\infty}f(x)\,\mathrm{d}x$同敛散.
\end{ti}

\begin{ti}
\textbf{定理7.13($Kummer$)}  设$\displaystyle\sum a_n$,$\displaystyle\sum b_n$为正项级数,如果$n$充分大时, \\
(1)$\dfrac{1}{b_n}\cdot\dfrac{a_n}{a_{n+1}}-\dfrac{1}{b_{n+1}}\geq\lambda>0$,则$\displaystyle\sum a_n$收敛; \\
(2)$\dfrac{1}{b_n}\cdot\dfrac{a_n}{a_{n+1}}-\dfrac{1}{b_{n+1}}\leq 0$且$\displaystyle\sum b_n$发散,则$\displaystyle\sum a_n$发散.
\end{ti}

\textbf{证明}:(1)条件可改写为$\dfrac{a_k}{b_k}-\dfrac{a_{k+1}}{b_{k+1}}\geq\lambda a_{k+1}$,这说明当$n\geq N$时,
$$S_{n+1}\leq S_N+\frac{1}{\lambda}\sum_{k=N}^{n}\left(\frac{a_k}{b_k}-\frac{a_{k+1}}{b_{k+1}}\right)\leq S_N+\frac{1}{\lambda}\cdot\frac{a_N}{b_N},$$
即$\{S_n\}$有上界,从而$\displaystyle\sum a_n$收敛. \\
(2)由已知条件可得$\dfrac{a_n}{b_n}\leq\dfrac{a_{n+1}}{b_{n+1}}$,即$\dfrac{a_n}{b_n}$关于$n$单调递增,从而$a_n\geq\dfrac{a_N}{b_N}b_n$,由比较判别法立得欲证结论.

\begin{ti}
\textbf{命题7.14($Sapagof$判别法)} 设正数数列$\{a_n\}$单调递减,则$\displaystyle\lim_{n\to\infty}a_n=0$的充分必要条件是正项级数$\displaystyle\sum_{n=1}^{\infty}\left(1-\dfrac{a_{n+1}}{a_n}\right)$发散.
\end{ti}

\textbf{注}:本命题有几种等价形式: \\
(1) 设$\{a_n\}$是单调递增的正数数列,则该数列与级数$\displaystyle\sum_{n=1}^{\infty}\left(1-\dfrac{a_n}{a_{n+1}}\right)$同敛散.\\
(2) 若正项级数$\displaystyle\sum_{n=1}^{\infty}a_n$的部分和数列为$\{S_n\}$,则$\displaystyle\sum_{n=1}^{\infty}a_n\text{与}\sum_{n=1}^{\infty}\dfrac{a_n}{S_n}$同敛散.

\begin{ti}
\textbf{命题7.15($\text{Du Bois-Reymond}$定理)} 对于一个给定的收敛正项级数$\displaystyle\sum_{n=1}^{\infty}a_n$,一定存在一个收敛正项级数$\displaystyle\sum_{n=1}^{\infty}b_n$,
使得$\displaystyle\lim_{n\to\infty}\dfrac{a_n}{b_n}=0$.
\end{ti}

\begin{ti}
\textbf{命题7.16($Carleman$不等式)} 设$\displaystyle\sum_{n=1}^{\infty}a_n$为收敛的正项级数,则成立
$$\sum_{n=1}^{\infty}\sqrt[n]{a_1a_2\cdots a_n}\leq e\sum_{n=1}^{\infty}a_n $$,
其中右边的系数$e$不能再改进.
\end{ti}

\begin{ti}
\textbf{定义7.17(无穷乘积)}  设$\{p_n\}$为一列实数,称形式乘积$\displaystyle\prod_{n=1}^{\infty}p_n=p_1p_2\cdots p_n\cdots$为\textbf{无穷乘积},记$P_n=p_1p_2\cdots p_n$,称为\textbf{部分乘积}.
如果数列$\{P_n\}$的极限存在,且极限为实数或正负无穷,则此极限称为无穷乘积的值,记为$\displaystyle\prod_{n=1}^{\infty}p_n=\lim_{n\to\infty}P_n$.当极限为非零实数时,称无穷乘积是\textbf{收敛}的,否则就称它是发散的.
\end{ti}

\begin{ti}
\textbf{命题7.18(无穷乘积的收敛判别法)}  设$p_n$均大于零,记$p_n=1+a_n$,则\\
(1)无穷乘积$\displaystyle\prod_{n=1}^{\infty}p_n$收敛当且仅当级数$\displaystyle\sum_{n=1}^{\infty}\ln p_n$收敛;\\
(2)如果$n$充分大时$a_n$不变号,则无穷乘积$\displaystyle\prod_{n=1}^{\infty}p_n$收敛当且仅当级数$\displaystyle\sum_{n=1}^{\infty}a_n$收敛;\\
(3)如果级数$\displaystyle\sum_{n=1}^{\infty}a_n$和$\displaystyle\sum_{n=1}^{\infty}a_n^2$均收敛,则无穷乘积$\displaystyle\prod_{n=1}^{\infty}p_n$也收敛.
\end{ti}


\begin{ti}
\textbf{定义7.19(级数的乘积)}  设级数$\displaystyle\sum_{n=0}^{\infty}a_n$与$\displaystyle\sum_{n=0}^{\infty}b_n$绝对收敛,定义它们的乘积为级数$\displaystyle\sum_{n=0}^{\infty}c_n$,其中$c_n=\displaystyle\sum_{i+j=n}a_i b_j$,$n\geq 0$.这种乘积称为\textbf{$Cauchy$乘积}.则乘积级数$\displaystyle\sum_{n=0}^{\infty}c_n$也绝对收敛,且
$$\sum_{n=0}^{\infty}c_n=\left(\sum_{n=0}^{\infty}a_n\right)\left(\sum_{n=0}^{\infty}b_n\right)$$.
\end{ti}

\begin{ti}
\textbf{定理7.20($Mertens$)}  如果$\displaystyle\sum_{n=0}^{\infty}a_n$和$\displaystyle\sum_{n=0}^{\infty}b_n$收敛,且至少其中一个级数绝对收敛,则它们的乘积级数也收敛,且
$$\sum_{n=0}^{\infty}c_n=\left(\sum_{n=0}^{\infty}a_n\right)\left(\sum_{n=0}^{\infty}b_n\right)$$.
\end{ti}

\textbf{证明}:不妨设$\displaystyle\sum a_n$绝对收敛.分别记$A_n=\displaystyle\sum_{k=0}^{n}a_k$,$B_n=\displaystyle\sum_{k=0}^{n}b_k$,$C_n=\displaystyle\sum_{k=0}^{n}c_k$.则$A_n\to A$,$B_n\to B$,而
$$C_n=\sum_{i+j\leq n}a_ib_j=a_0B_n+a_1B_{n-1}+\cdots+a_nB_0=A_nB+\delta_n,$$
其中$\delta_n=a_0(B_n-B)+a_1(B_{n-1}-B)+\cdots+a_n(B_0-B)$.我们只要证明$\delta_n\to 0$即可.因为$B_n\to B$,故$\{B_n\}$关于$n$有界,从而存在$K$,使得$|B_n-B|\leq K$对一切$n\geq 0$成立.由于$\displaystyle\sum a_n$绝对收敛,故任给$\varepsilon>0$,存在$N_0$,当$n>N_0$时,$|a_{N_0+1}|+|a_{N_0+2}|+\cdots+|a_n|<\dfrac{\varepsilon}{2K+1}$.记$L=|a_0|+|a_1|+\cdots+|a_{N_0}|$.由于$B_n-B\to 0$,故存在$N_1$,当$n>N_1$时,$|B_n-B|<\dfrac{\varepsilon}{2L+1}$.从而当$n>N_0+N_1$时,有
$$|\delta_n|\leq\sum_{k=0}^{N_0}|a_k||B_{n-k}-B|+(|a_{N_0+1}|+|a_{N_0+2}|+\cdots+|a_n|)K<\frac{\varepsilon}{2L+1}L+\frac{\varepsilon}{2K+1}K<\varepsilon,$$
这说明$\delta_n\to 0$,因而$C_n=A_nB+\delta_n\to AB$.

\begin{ti}
\textbf{引理7.21($Abel$)}  设级数$\displaystyle\sum_{n=0}^{\infty}C_n$收敛,则$\displaystyle\lim_{x\to 1^-}\sum_{n=0}^{\infty}C_n x^n=\sum_{n=0}^{\infty}C_n$.
\end{ti}

\begin{ti}
\textbf{定理7.22($Abel$)}  设级数$\displaystyle\sum_{n=0}^{\infty}a_n$,$\displaystyle\sum_{n=0}^{\infty}b_n$以及它们的乘积$\displaystyle\sum_{n=0}^{\infty}c_n$均收敛,则
$$\sum_{n=0}^{\infty}c_n=\left(\sum_{n=0}^{\infty}a_n\right)\left(\sum_{n=0}^{\infty}b_n\right)$$.
\end{ti}

\textbf{证明}:当$x\in[0,1)$时,级数$\displaystyle\sum_{n=0}^{\infty}a_n x^n$和$\displaystyle\sum_{n=0}^{\infty}b_n x^n$绝对收敛,它们的乘积级数为$\displaystyle\sum_{n=0}^{\infty}c_n x^n$.根据$Cauchy$定理,有
$$\sum_{n=0}^{\infty}c_n x^n=\left(\sum_{n=0}^{\infty}a_n x^n\right)\left(\sum_{n=0}^{\infty}b_n x^n\right).$$
令$x\to 1^-$,由上述$Abel$引理即得欲证结论.

\begin{ti}
\textbf{定理7.23(重排定理)}  绝对收敛级数的项可以任意重排而不改变其和.
\end{ti}

\textbf{证明}:设级数$\displaystyle\sum a_n$绝对收敛.如果$a_n$为正项级数,则由基本判别法不难看出,将它的各项重新排列后不会影响其敛散性,如果收敛的话重排也不改变级数的和.对于一般的级数,由等式
$$a_n=\dfrac{a_n+|a_n|}{2}-\dfrac{|a_n|-a_n}{2}$$
以及正项级数$\displaystyle\sum\dfrac{a_n+|a_n|}{2}$与$\displaystyle\sum\dfrac{|a_n|-a_n}{2}$的收敛性可知,重排后的级数也绝对收敛,且其和不变.

\begin{ti}
\textbf{定理7.24($Riemann$重排定理)}  如果级数$\displaystyle\sum a_n$条件收敛,则可以将它重排为一个收敛级数,使得重排后的级数和为任意指定的实数.
\end{ti}

\textbf{证明}:设$\xi\in\mathbb{R}$,我们将找到$\displaystyle\sum a_n$的一个重排,使得它的和为$\xi$.首先我们注意到,如果$\displaystyle\sum a_n$条件收敛,则$\displaystyle\sum|a_n|$发散,由$|a_n|=a_n^++a_n^-$,$a_n=a_n^+-a_n^-$即知,级数$\displaystyle\sum a_n^+$和$\displaystyle\sum a_n^-$均发散到$+\infty$,这两个级数的部分和分别记为$S_n^+$和$S_n^-$.因为当$n$充分大时$S_n^+>\xi$,故可取最小的正整数$m_1$,使得$S_{m_1}^+>\xi$,此时成立$\xi\geq S_{m_1-1}^+$.因为$n$充分大时,$S_{m_1}^+-S_n^-<\xi$,故可取最小的正整数$n_1$,使得$S_{m_1}^+-S_{n_1}^-<\xi$,同理有$\xi\leq S_{m_1}^+-S_{n_1-1}^-$.下面再取最小的正整数$m_2$,使得$S_{m_2}^+-S_{n_1}^->\xi\geq S_{m_2-1}^+-S_{n_1}^-$,以及最小的正整数$n_2$,使得$S_{m_2}^+-S_{n_2}^-<\xi\leq S_{m_2}^+-S_{n_2-1}^-$.如此继续下去,我们得到递增数列$m_1<m_2<\cdots$和$n_1<n_2<\cdots$,使得
$$S_{m_k}^+-S_{n_{k-1}}^->\xi\geq S_{m_k-1}^+-S_{n_{k-1}}^-,\qquad S_{m_k}^+-S_{n_k}^-<\xi\leq S_{m_k}^+-S_{n_k-1}^-$$
对任意$k$均成立.由$a_n\to 0$即知,下面的级数
$$a_1^++\cdots+a_{m_1}^+-a_1^--\cdots-a_{n_1}^-+a_{m_1+1}^++\cdots+a_{m_2}^+-a_{n_1+1}^--\cdots-a_{n_2}^-+\cdots$$
收敛到$\xi$.注意$a_n^+$和$a_n^-$在这个级数中都依次出现了,因此它可以看成是原级数$\displaystyle\sum a_n$的一个重排.

\subsection{函数项级数}
\begin{ti}
\textbf{定义7.25(一致收敛)}  设$I$为区间,$\{f_n\}$为$I$中定义的一列函数.如果对任意的$\varepsilon>0$,均存在与$x$无关的正整数$N=N(\varepsilon)$,使得当$n>N,x\in I$时,$\left\lvert f_n(x)-f(x)\right\rvert<\varepsilon$总成立,则称函数列$\{f_n\}$在$I$中\textbf{一致收敛}于$f$,记为$f_n\Rightarrow f$.
对于函数项级数$\displaystyle\sum f_n(x)$,若其部分和函数列一致收敛,则称该级数在$I$上\textbf{一致收敛}.
\end{ti}

\begin{ti}
\textbf{定理7.26(一致收敛的$Cauchy$准则)}  定义在$I$中的函数列$\{f_n\}$一致收敛当且仅当对任意的$\varepsilon>0$,存在$N=N(\varepsilon)$,使得当$m,n>N,x\in I$时,$\left\lvert f_m(x)-f_n(x)\right\rvert<\varepsilon$.
\end{ti}

\begin{ti}
\textbf{定理7.27(一致收敛与连续性)}  设$\{f_n\}$在区间$I$中一致收敛于$f$.如果每一个$f_n$均为连续函数,则$f$也是连续函数.
\end{ti}

\begin{ti}
\textbf{定理7.28(极限交换)}  设$\{f_n\}$在$x_0$的一个空心邻域中一致收敛于$f$.如果$f_n$在$x_0$处的函数极限为$a_n$,则数列极限$\displaystyle\lim_{n\to\infty}a_n$以及函数极限$\displaystyle\lim_{x\to x_0}f(x)$均存在,且这两个极限相等,即
$$\lim_{n\to\infty}\lim_{x\to x_0}f_n(x)=\lim_{x\to x_0}\lim_{n\to\infty}f_n(x).$$
\end{ti}

\textbf{证明}:由一致收敛性,任给$\varepsilon>0$,存在$N_0$,当$n>N_0$时,在$x_0$的空心邻域中均有$|f_n(x)-f(x)|<\varepsilon$.特别地,当$m,n>N_0$时,$|f_m(x)-f_n(x)|\leq|f_m(x)-f(x)|+|f(x)-f_n(x)|<2\varepsilon$.在上式中令$x\to x_0$,得$|a_m-a_n|\leq 2\varepsilon$对一切$m,n>N_0$成立.由$Cauchy$准则可知$\{a_n\}$收敛,其极限记为$\alpha$.下证$\displaystyle\lim_{x\to x_0}f(x)=\alpha$.任给$\varepsilon>0$,由刚才的证明,存在$N$,使得$|a_N-\alpha|<\dfrac{\varepsilon}{3}$,$|f(x)-f_N(x)|<\dfrac{\varepsilon}{3}$.因为$\displaystyle\lim_{x\to x_0}f_N(x)=a_N$,故存在$\delta>0$,当$0<|x-x_0|<\delta$时,$|f_N(x)-a_N|<\dfrac{\varepsilon}{3}$.因此,当$0<|x-x_0|<\delta$时,
$$|f(x)-\alpha|\leq|f(x)-f_N(x)|+|f_N(x)-a_N|+|a_N-\alpha|<\frac{\varepsilon}{3}+\frac{\varepsilon}{3}+\frac{\varepsilon}{3}=\varepsilon,$$
即$\displaystyle\lim_{x\to x_0}f(x)=\alpha$.

\begin{ti}
\textbf{定义7.29(一致有界)} 设在区间$I$上给定一个函数序列
$$f_1(x),f_2(x),...,f_n(x),...,$$
如果存在常数$K>0$,使得不等式
$$\left\lvert f_n(x)\right\rvert\leq K\quad(x\in I)$$
对一切$n=1,2,...$都成立,则称这个函数序列在区间$I$上一致有界.
\end{ti}

\begin{ti}
\textbf{定义7.30(等度连续)} 设在区间$I$上给定一个函数序列
$$f_1(x),f_2(x),...,f_n(x),...,$$
如果对任意的正数$\varepsilon$,存在正数$\delta=\delta(\varepsilon)$,使得只要$x_1,x_2\in I$和$\left\lvert x_1-x_2\right\rvert<\delta$,就有
$$\left\lvert f_n(x_1)-f_n(x_2)\right\rvert<\varepsilon\quad(n=1,2,...),$$
则称这个函数序列在区间$I$上是等度连续的.
\end{ti}

\begin{ti}
\textbf{定理7.31($Ascoli$引理)} 设函数序列$f_1(x),f_2(x),...,f_n(x),...$在有限闭区间$[a,b]$上是一致有界和等度连续的,则可以选取它的一个子序列
$$f_{n_1}(x),f_{n_2}(x),...,f_{n_k}(x),...,$$
使该子序列在区间$I$上是一致收敛的.
\end{ti}

\textbf{证明}:将$[a,b]$中所有有理数排成一列(由有理数的可数性可知这样做是合法的):
$$r_1,r_2,...$$
因为$\{f_n\}$一致有界,所以对任意固定的$r_j$,数列$\{f_n(r_j)\}$是有界数列.由$Bolzano$定理,它有收敛子列.
现在用对角线法:
先取子列$\{f_{1,k}\}$,使得$\{f_{1,k}(r_1)\}$收敛;
再在$\{f_{1,k}\}$中取子列$\{f_{2,k}\}$,使得$\{f_{2,k}(r_2)\}$收敛,同时因为他是$\{f_{1,k}\}$的子列,所以在$r_1$处也仍然收敛;
继续这样做下去,得到一串嵌套的子列:
$$\{f_{1,k}\}\supset\{f_{2,k}\}\supset\{f_{3,k}\}\supset ...$$
使得对每个有理点$r_j$,子列$\{f_{j,k}\}$在$r_j$处收敛.
最后取对角线上的函数:$g_k=f_{k,k}$. 
于是对任意固定的有理点$r_j$,当$k\geq j$时,$g_k(r_j)$收敛.\\

任给$\varepsilon>0$,因为原序列$\{f_n\}$等度连续,所以它的子列$\{g_k\}$也等度连续.因此存在$\delta>0$,使得只要$\left\lvert x-y\right\rvert<\delta$,
则对所有的$k$都有$\left\lvert g_k(x)-g_k(y)\right\rvert<\dfrac{\varepsilon}{3}$.\\\
现在将$[a,b]$等分成有限段,使每一段长度小于$\delta$,
例如取$N$足够大,使得$\dfrac{b-a}{N}<\delta$.在每一小段中取一个有理点$q_1,q_2,...,q_N$.由有理数的稠密性,这是可以做到的.
由于$\{g_k\}$在每一个$q_i$处均收敛,而这里只有有限个点$q_1,q_2,...,q_N$,由$Cauchy$准则,存在$K$,使得当$k,l\geq K$时,对所有$i=1,...,N$,都有
$$\left\lvert g_k(q_i)-g_l(q_i)\right\rvert<\dfrac{\varepsilon}{3}.$$
现在任取$x\in[a,b]$,它一定属于某个小段,设在这个小段中取到的有理点为$q_i$,则
$$\left\lvert x-q_i\right\rvert<\delta.$$
于是当$k,l\geq K$时,
$$\left\lvert g_k(x)-g_l(x)\right\rvert\leq\left\lvert g_k(x)-g_k(q_i)\right\rvert$$
$$\phantom{\left\lvert g_k(x)-g_l(x)\right\rvert\leq}+\left\lvert g_k(q_i)-g_l(q_i)\right\rvert$$
$$\phantom{\left\lvert g_k(x)-g_l(x)\right\rvert\leq}+\left\lvert g_l(q_i)-g_l(x)\right\rvert .$$
从而有
$$\left\lvert g_k(x)-g_l(x)\right\rvert<\varepsilon.$$
这对任意$x\in[a,b]$都成立,所以$\{g_k\}$在$[a,b]$上一致收敛.得证.

\begin{ti}
\textbf{定理7.32($Dini$)}  设$\{f_n\}$为$[a,b]$中的一列非负连续函数,且对每个$x\in[a,b]$,$\{f_n(x)\}$关于$n$单调递减趋于$0$,则函数列$\{f_n\}$一致收敛于$0$.
\end{ti}

\textbf{证明}:任给$\varepsilon>0$,我们要证明存在$N$,使得当$n>N,x\in[a,b]$时,$0\leq f_n(x)<\varepsilon$.即要证$n$充分大以后$A_n=\{x\in[a,b]\;|\;f_n(x)\geq\varepsilon\}$为空集.因为$f_n$关于$n$单调递减,因此
$$A_1\supset A_2\supset\cdots\supset A_n\supset A_{n+1}\supset\cdots,$$
这说明我们只要证明某一个$A_n$是空集即可.\\
(反证法)假设$A_n$均为非空集,取$x_n\in A_n$,则$\{x_n\}$为$[a,b]$中的有界点列,根据$Bolzano$定理,它存在收敛子列$\{x_{n_i}\}$,设此子列收敛到$x_0\in[a,b]$.由上面的包含关系可知$A_k\supset A_{n_k}\supset\{x_{n_k},x_{n_k+1},\cdots\}$.因为$f_k$在$x_0$处连续,我们有
$$f_k(x_0)=\lim_{i\to\infty}f_k(x_{n_i})\geq\varepsilon.$$
上式对每一个$k\geq 1$均成立,但这和$\{f_n(x_0)\}$收敛于零相矛盾.

\begin{ti}
\textbf{推论7.33}  设$f_n(x)$为非负连续函数,如果函数项级数$\displaystyle\sum_{n=1}^{\infty}f_n(x)$在闭区间中收敛于连续函数$f$,则必一致收敛于$f$.
\end{ti}

\begin{ti}
\textbf{定理7.34(逐项积分)}  设$f_n\in R[a,b]$,如果$\{f_n\}$一致收敛于$f$,则$f\in R[a,b]$,且
$$\lim_{n\to\infty}\int_a^b f_n(x)\,\mathrm{d}x=\int_a^b \lim_{n\to\infty}f_n(x)\,\mathrm{d}x=\int_a^b f(x)\,\mathrm{d}x$$.
\end{ti}

\textbf{证明}:先来说明$f$可积.任给$\varepsilon>0$,存在$N=N(\varepsilon)$,当$n\geq N,x\in[a,b]$时,
$|f_n(x)-f(x)|\leq\dfrac{\varepsilon}{4(b-a)}$.因为$f_N$是可积函数,故存在$[a,b]$的分割$\pi$,使得$\displaystyle\sum_i\omega_i(f_N)\Delta x_i<\dfrac{\varepsilon}{2}$.对于每一个小区间$[x_{i-1},x_i]$,由上式可得
$$\omega_i(f)\leq\omega_i(f_N)+\dfrac{\varepsilon}{2(b-a)}.$$
因此
$$\sum_i\omega_i(f)\Delta x_i\leq\sum_i\omega_i(f_N)\Delta x_i+\dfrac{\varepsilon}{2(b-a)}(b-a)\leq\dfrac{\varepsilon}{2}+\dfrac{\varepsilon}{2}=\varepsilon,$$
由可积函数的充要条件即知$f\in R[a,b]$.进一步,当$n\geq N$时,我们有
$$\left\lvert\int_a^b f_n(x)\,\mathrm{d}x-\int_a^b f(x)\,\mathrm{d}x\right\rvert\leq\int_a^b|f_n(x)-f(x)|\,\mathrm{d}x\leq\dfrac{\varepsilon}{4},$$
这说明$\displaystyle\lim_{n\to\infty}\int_a^b f_n(x)\,\mathrm{d}x=\int_a^b f(x)\,\mathrm{d}x$.

\begin{ti}
\textbf{推论7.35}  设$\displaystyle\sum_{n=1}^{\infty}f_n(x)$在$[a,b]$中一致收敛于$f$.如果$f_n$均为可积函数,则$f$也是可积函数,且
$$\sum_{n=1}^{\infty}\int_a^b f_n(x)\,\mathrm{d}x=\int_a^b\sum_{n=1}^{\infty}f_n(x)\,\mathrm{d}x=\int_a^b f(x)\,\mathrm{d}x.$$
\end{ti}

\begin{ti}
\textbf{定理7.36(逐项求导)}  设$\{f_n\}$为$[a,b]$中的一列可导函数,$c\in[a,b]$.如果$\{f_n(c)\}$收敛,$\{f'_n\}$一致收敛于$g$,
则$\{f_n\}$一致收敛于可导函数$f$,且$f'=g$.
\end{ti}

\textbf{证明}:由$\{f'_n\}$一致收敛可知,任给$\varepsilon>0$,存在$N$,当$m,n>N,x\in[a,b]$时,$|f'_m(x)-f'_n(x)|<\varepsilon$.此时,由微分中值定理可得
$$|[f_m(x)-f_m(c)]-[f_n(x)-f_n(c)]|=|[f'_m(\xi)-f'_n(\xi)](x-c)|\leq(b-a)\varepsilon,$$
这说明$\{f_n-f_n(c)\}$一致收敛,从而$\{f_n\}$一致收敛到某个函数$f(x)$.\\
我们要说明$f$可导.为此,任取$x_0\in[a,b]$,令
$$g_n(x)=\begin{cases}\dfrac{f_n(x)-f_n(x_0)}{x-x_0},&x\neq x_0,\\[2mm] f'_n(x_0),&x=x_0.\end{cases}$$
由$f_n$可导可知$g_n$为连续函数.类似于刚才的论证,由微分中值定理,有$|g_m(x)-g_n(x)|=|f'_m(\xi)-f'_n(\xi)|\leq\varepsilon$,这说明$\{g_n\}$一致收敛,其极限为
$$\widehat{g}(x)=\begin{cases}\dfrac{f(x)-f(x_0)}{x-x_0},&x\neq x_0,\\[2mm] g(x_0),&x=x_0.\end{cases}$$
根据前节定理,$\;\widehat{g}$为连续函数.特别地,$f$在$x_0$处可导且导数为$g(x_0)$.

\begin{ti}
\textbf{推论7.37}  设$\{f_n\}$为$[a,b]$中一列可导函数,$c\in[a,b]$.如果$\displaystyle\sum_{n=1}^{\infty}f_n(c)$收敛,$\displaystyle\sum_{n=1}^{\infty}f'_n(x)$一致收敛,则$\displaystyle\sum_{n=1}^{\infty}f_n(x)$一致收敛,其和函数可导,且
$$\left(\sum_{n=1}^{\infty}f_n(x)\right)'=\sum_{n=1}^{\infty}f'_n(x).$$
\end{ti}

\begin{ti}
\textbf{引理7.38($Abel$)}  设幂级数$\displaystyle\sum_{n=0}^{\infty}a_n x^n$在$x=x_1\,(x_1\neq 0)$处收敛,则它在区间$|x|<|x_1|$中绝对收敛.因此,幂级数在$x=x_2$处发散意味着在$|x|>|x_2|$中均发散.
\end{ti}

\textbf{证明}:设$\displaystyle\sum_{n=0}^{\infty}a_n x_1^n$收敛,则存在$M>0$,使得$|a_n x_1^n|\leq M$均成立.这说明
$$\sum_{n=0}^{\infty}|a_n x^n|=\sum_{n=0}^{\infty}|a_n x_1^n|\left|\frac{x}{x_1}\right|^n\leq M\sum_{n=0}^{\infty}\left|\frac{x}{x_1}\right|^n,$$
因此当$|x|<|x_1|$时,$\displaystyle\sum_{n=0}^{\infty}a_n x^n$绝对收敛.

\begin{ti}
\textbf{定理7.39($Cauchy\text{-}Hadamard$)}  对幂级数$\displaystyle\sum_{n=0}^{\infty}a_n x^n$,记$\rho=\displaystyle\limsup_{n\to\infty}\sqrt[n]{\left\lvert a_n\right\rvert}$,则\\
(1)$\rho=0$时,幂级数在$(-\infty,\infty)$中绝对收敛;\\
(2)$\rho=+\infty$时,幂级数仅在$x=0$处收敛;\\
(3)$0<\rho<\infty$时,幂级数在$\left\lvert x\right\rvert<\dfrac{1}{\rho}$中绝对收敛,在$\left\lvert x\right\rvert>\dfrac{1}{\rho}$之外发散.此时,称$\dfrac{1}{\rho}$为\textbf{收敛半径}.
\end{ti}


\begin{ti}
\textbf{定理7.40($Abel$连续性定理)}  设幂级数$\displaystyle\sum_{n=0}^{\infty}a_n x^n$的收敛半径为$R\,(0<R<\infty)$.如果$\displaystyle\sum_{n=0}^{\infty}a_n R^n$收敛,则
$$\lim_{x\to R^-}\sum_{n=0}^{\infty}a_n x^n=\sum_{n=0}^{\infty}a_n R^n;$$
如果$\displaystyle\sum_{n=0}^{\infty}a_n(-R)^n$收敛,则
$$\lim_{x\to -R^+}\sum_{n=0}^{\infty}a_n x^n=\sum_{n=0}^{\infty}a_n(-R)^n.$$
\end{ti}

\textbf{证明}:设$\displaystyle\sum a_nR^n$收敛.当$x\in[0,R]$时,$\displaystyle\sum a_nx^n=\sum a_nR^n\left(\dfrac{x}{R}\right)^n$,其中$\left(\dfrac{x}{R}\right)^n$关于$n$单调且$\leq 1$,由数项级数的$Abel$判别法,$\displaystyle\sum a_nx^n$在$[0,R]$上一致收敛,故$\displaystyle\lim_{x\to R^-}\sum a_nx^n=\sum a_nR^n$.对$x\in[-R,0]$的情形同理,即得另一等式.

\section{\texorpdfstring{$Fourier$}{Fourier}分析}
\begin{ti}
\textbf{定义8.1($Fourier$系数与$Fourier$级数)}  设$f$在$[-\pi,\pi]$中$Riemann$可积或广义绝对可积.令
$$a_0=\dfrac{1}{\pi}\int_{-\pi}^{\pi}f(x)\,\mathrm{d}x,a_n=\dfrac{1}{\pi}\int_{-\pi}^{\pi}f(x)\cos(nx)\,\mathrm{d}x,$$
$$b_n=\dfrac{1}{\pi}\int_{-\pi}^{\pi}f(x)\sin(nx)\,\mathrm{d}x,\quad n=1,2,\ldots,$$
$a_0,a_n,b_n$称为$f$的\textbf{$Fourier$系数},形式和
$$\dfrac{a_0}{2}+\sum_{n=1}^{\infty}\bigl(a_n\cos(nx)+b_n\sin(nx)\bigr)$$
称为$f$的\textbf{$Fourier$级数}或$Fourier$展开,记为
$$\displaystyle f(x)\sim\dfrac{a_0}{2}+\sum_{n=1}^{\infty}\bigl(a_n\cos(nx)+b_n\sin(nx)\bigr).$$
\end{ti}


\begin{ti}
\textbf{命题8.2}  设$f$在$[-\pi,\pi]$中可积或广义绝对可积,则其$Fourier$系数形如
$$a_n=o(1),\quad b_n=o(1)\quad(n\to\infty).$$
\end{ti}

\begin{ti}
\textbf{命题8.3}  设$f\in C^0[-\pi,\pi]$,$f(-\pi)=f(\pi)$.如果$f$是分段线性函数,则其$Fourier$系数形如
$$a_n=O\left(\frac{1}{n^2}\right),\quad b_n=O\left(\frac{1}{n^2}\right)\quad(n\to\infty).$$
\end{ti}

\begin{ti}
\textbf{定理8.4($Bessel$不等式)}  设$f\in R[-\pi,\pi]$,则
$$\dfrac{a_0^2}{2}+\sum_{n=1}^{\infty}\bigl(a_n^2+b_n^2\bigr)\leq\dfrac{1}{\pi}\int_{-\pi}^{\pi}f(x)^2\,\mathrm{d}x$$.
\end{ti}

\textbf{证明}:任取$g\in R[-\pi,\pi]$,设其$Fourier$系数为$\{\alpha_n,\beta_n\}$,记$S_n(x)=\dfrac{a_0}{2}+\displaystyle\sum_{k=1}^{n}(a_k\cos kx+b_k\sin kx)$为$f$的$Fourier$级数的部分和,则
$$\dfrac{1}{\pi}\int_{-\pi}^{\pi}S_n(x)g(x)\,\mathrm{d}x=\dfrac{1}{\pi}\left[\dfrac{a_0}{2}\int_{-\pi}^{\pi}g(x)\,\mathrm{d}x\right.$$
$$\phantom{\dfrac{1}{\pi}\int_{-\pi}^{\pi}S_n(x)g(x)\,\mathrm{d}x=}\left.+\sum_{k=1}^{n}\left(a_k\int_{-\pi}^{\pi}g(x)\cos kx\,\mathrm{d}x+b_k\int_{-\pi}^{\pi}g(x)\sin kx\,\mathrm{d}x\right)\right]$$
$$=\dfrac{1}{2}a_0\alpha_0+\sum_{k=1}^{n}(a_k\alpha_k+b_k\beta_k).$$
特别地,在上式中分别取$g(x)=S_n(x),f(x)$,注意到$S_n(x)$的$Fourier$展开是它自身,可以得到
$$\dfrac{1}{\pi}\int_{-\pi}^{\pi}S_n(x)^2\,\mathrm{d}x=\dfrac{1}{2}a_0^2+\sum_{k=1}^{n}(a_k^2+b_k^2),$$
$$\dfrac{1}{\pi}\int_{-\pi}^{\pi}S_n(x)f(x)\,\mathrm{d}x=\dfrac{1}{2}a_0^2+\sum_{k=1}^{n}(a_k^2+b_k^2),$$
由$\int_{-\pi}^{\pi}(f-S_n)^2\,\mathrm{d}x\geq0$得$\dfrac{a_0^2}{2}+\sum_{k=1}^n(a_k^2+b_k^2)\leq\dfrac{1}{\pi}\int_{-\pi}^{\pi}f^2\,\mathrm{d}x$.令$n\to\infty$即得结论.

\begin{ti}
\textbf{定理8.5($Dini$判别法)}  设$f$如前.如果存在$\delta>0$,使得$f$在$x$处的左极限$f(x-0)$和右极限$f(x+0)$均存在,且积分
$$\int_0^{\delta}\frac{|f(x+u)-f(x+0)|}{u}\,\mathrm{d}u,\qquad \int_0^{\delta}\frac{|f(x-u)-f(x-0)|}{u}\,\mathrm{d}u$$
绝对收敛,则$f$的$Fourier$展开在$x$处收敛于$\dfrac{f(x+0)+f(x-0)}{2}$.
\end{ti}

\textbf{证明}:其证明基本上是应用$Riemann-Lebesgue$引理.

\begin{ti}
\textbf{定理8.6($Dirichlet$收敛定理)}  设$f$以$2\pi$为周期,在一个周期内仅有有限个分段点,且每段内为$C^1$函数,$f$及$f'$在各段端点处均有有限的单侧极限.则对每个$x\in[-\pi,\pi]$,$f$的$Fourier$展开在$x$处收敛到
$$\dfrac{f(x+0)+f(x-0)}{2}$$
特别地,在$f$的连续点处,$Fourier$级数收敛到$f(x)$.
\end{ti}

\textbf{证明}:由$Dirichlet$核的计算,$Fourier$级数的部分和可写为
$$S_n(x)=\dfrac{1}{\pi}\int_0^{\pi}\bigl[f(x+u)+f(x-u)\bigr]\cdot\dfrac{\sin\left(n+\dfrac{1}{2}\right)u}{2\sin\dfrac{u}{2}}\,\mathrm{d}u.$$
记$D_n(u)=\dfrac{\sin\left(n+\frac12\right)u}{2\sin(u/2)}$.由$\displaystyle\int_0^\pi D_n(u)\,\mathrm{d}u=\dfrac{\pi}{2}$,上式可以改写为
$$S_n(x)=\dfrac{f(x+0)+f(x-0)}{2}+E_n(x),$$
其中
$$E_n(x)=\dfrac{1}{\pi}\int_0^{\pi}\dfrac{f(x+u)-f(x+0)}{2\sin\dfrac{u}{2}}\sin\left(n+\dfrac{1}{2}\right)u\,\mathrm{d}u$$
$$\phantom{E_n(x)=}+\dfrac{1}{\pi}\int_0^{\pi}\dfrac{f(x-u)-f(x-0)}{2\sin\dfrac{u}{2}}\sin\left(n+\dfrac{1}{2}\right)u\,\mathrm{d}u.$$
由所设的分段$C^1$及端点单侧导数有限条件,$f(x\pm u)-f(x\pm0)=O(u)$($u\to0+$),故上述两个差商在$u=0$附近有界,在$[0,\pi]$上分段连续.应用$Riemann-Lebesgue$引理即得$E_n(x)\to0$.在$f$的连续点处,$f(x+0)=f(x-0)=f(x)$,故收敛到$f(x)$.

\begin{ti}
\textbf{定理8.7($Weierstrass$逼近定理)}  设$f\in C^0[-\pi,\pi]$,$f(-\pi)=f(\pi)$.则任给$\varepsilon>0$,存在三角多项式$T(x)$,使得$|f(x)-T(x)|<\varepsilon$在$[-\pi,\pi]$中处处成立.
\end{ti}

\textbf{证明}:首先,连续函数$f$可以用分段线性函数一致逼近,即存在周期为$2\pi$的连续分段线性函数列$\{f_n\}$,使得$f_n$在$[-\pi,\pi]$中一致收敛于$f$.其次,由命题8.3的系数估计与$Weierstrass$判别法,各$f_n$的$Fourier$展开一致收敛;再由$Dirichlet$收敛定理,其和为$f_n$.取$n$充分大,使$|f(x)-f_n(x)|<\dfrac{\varepsilon}{2}$处处成立;再取$f_n$的$Fourier$展开的充分多项的部分和$T$,使$|f_n(x)-T(x)|<\dfrac{\varepsilon}{2}$处处成立,则$|f(x)-T(x)|<\varepsilon$在$[-\pi,\pi]$中处处成立.

\begin{ti}
\textbf{定理8.8(一致收敛)}  设$f$是以$2\pi$为周期的连续函数,如果$f$在$[-\pi,\pi]$上可以写成两个单调函数之差(即在一个周期内有界变差),则其$Fourier$级数一致收敛于$f$自身.特别地,$Lipschitz$周期函数的$Fourier$级数一致收敛于自身.
\end{ti}

\textbf{证明}:记$D_n(u)=\dfrac{\sin(n+\frac12)u}{2\sin(u/2)}$.由核的积分式和$\int_0^\pi D_n(u)\,\mathrm{d}u=\pi/2$可得
$$S_n(f)(x)-f(x)=\frac{1}{\pi}\int_0^\pi h_x(u)D_n(u)\,\mathrm{d}u,\quad h_x(u)=f(x+u)+f(x-u)-2f(x).$$
先注意存在与$n$无关的常数$C$,使$\left|\int_0^vD_n(u)\,\mathrm{d}u\right|\leq C$对所有$0\leq v\leq\pi$成立.事实上,写$D_n(u)=\dfrac{\sin(n+\frac12)u}{u}\dfrac{u}{2\sin(u/2)}$;后一因子在$[0,\pi]$上连续可微,而前一因子的原函数由在$u=(n+\frac12)^{-1}$处分段估计并分部积分可得一致有界,再作一次分部积分即得此界.\\
记$\operatorname{Var}_I f$为$f$在区间$I$上的总变差,$V=\operatorname{Var}_{[-\pi,\pi]}f$.因为$f$连续且有界变差,累积变差函数的跳跃量等于$f$在该点的跳跃量,故它连续.把$f$按周期延拓,并记$\omega_V(\delta)=\sup_{|I|\leq\delta}\operatorname{Var}_I f$,则累积变差函数的一致连续性给出$\omega_V(\delta)\to0$($\delta\to0+$).对$h_x(0)=0$在$[0,\delta]$上应用有限求和分部公式并取极限,由前述核原函数的统一界得到
$$\left|\int_0^\delta h_x(u)D_n(u)\,\mathrm{d}u\right|\leq C\operatorname{Var}_{[0,\delta]}h_x\leq2C\omega_V(\delta),$$
右端与$x,n$无关.固定$\delta>0$后,令$q_x(u)=h_x(u)/(2\sin(u/2))$.在$[\delta,\pi]$上,分母的倒数有界且有界变差,而$\lVert h_x\rVert_\infty\leq4\lVert f\rVert_\infty$,$\operatorname{Var}_{[\delta,\pi]}h_x\leq2V$,故$\lVert q_x\rVert_\infty+\operatorname{Var}_{[\delta,\pi]}q_x\leq C_\delta$,其中$C_\delta$与$x,n$无关.再次用有限求和分部公式及$\left|\int_s^t\sin(n+\frac12)u\,\mathrm{d}u\right|\leq2/(n+\frac12)$,得到
$$\sup_x\left|\int_\delta^\pi h_x(u)D_n(u)\,\mathrm{d}u\right|\leq\frac{C_\delta}{n+\frac12}\longrightarrow0.$$
先令$n\to\infty$,再令$\delta\to0+$,即得$\sup_x|S_n(f)(x)-f(x)|\to0$.最后,若$f$的$Lipschitz$常数为$L$,则在$[-\pi,\pi]$上$f(t)=[f(t)+Lt]-Lt$为两个单调函数之差,故适用上述结论.

\begin{ti}
\textbf{定理8.9($\text{Fejér}$)}  设$f$连续且以$2\pi$为周期,则$\{\sigma_n(f)\}$一致收敛于$f$,其中$\sigma_n(f)$为$f$的$Fourier$级数部分和的算术平均:
$$\sigma_n(f)(x)=\frac{1}{n+1}\sum_{k=0}^{n}S_k(x),$$
$S_k(x)$为$f$的$Fourier$级数的第$k$个部分和.
\end{ti}

\textbf{证明}:$\sigma_n(f)(x)$可以表示为
$$\sigma_n(f)(x)=\frac{1}{2\pi}\int_{-\pi}^{\pi}F_n(u)f(x+u)\,\mathrm{d}u,$$
其中$F_n(u)$为$Fej\acute{e}r$核,它满足如下的性质:\\
(1)$F_n(u)\geq 0$;\\
(2)$\dfrac{1}{2\pi}\displaystyle\int_{-\pi}^{\pi}F_n(u)\,\mathrm{d}u=1$;\\
(3)对任意的$\delta>0$,$\displaystyle\lim_{n\to\infty}\frac{1}{2\pi}\int_{\delta\leq|u|\leq\pi}F_n(u)\,\mathrm{d}u=0$.\\
由题设可知$f$一致连续.任给$\varepsilon>0$,存在$\delta\in(0,\pi)$,当$|u|<\delta,x\in\mathbb{R}$时$|f(x+u)-f(x)|<\dfrac{\varepsilon}{2}$.记$M=\max|f|$,利用$Fej\acute{e}r$核的性质可得
$$|\sigma_n(f)(x)-f(x)|\leq\frac{1}{2\pi}\int_{-\pi}^{\pi}F_n(u)|f(x+u)-f(x)|\,\mathrm{d}u<\frac{\varepsilon}{2}+4M\cdot\frac{1}{2\pi}\int_{\delta\leq|u|\leq\pi}F_n(u)\,\mathrm{d}u,$$
当$n$充分大时,上式右端第二项小于$\dfrac{\varepsilon}{2}$,故$|\sigma_n(f)(x)-f(x)|<\varepsilon$对一切$x$成立,即$\{\sigma_n(f)\}$一致收敛于$f$.

\begin{ti}
\textbf{推论8.10($Dirichlet\text{-}Jordan$)}  设$f$可积,如果$f$在区间$[x_0-h,x_0+h]$中是单调函数,则其$Fourier$级数在$x_0$处收敛于
$$\frac{f(x_0+0)+f(x_0-0)}{2}.$$
\end{ti}

\begin{ti}
\textbf{定理8.11($Parseval$恒等式)}  设$f\in R[-\pi,\pi]$,其$Fourier$展开为$f(x)\sim\dfrac{a_0}{2}+\displaystyle\sum_{n=1}^{\infty}\bigl(a_n\cos(nx)+b_n\sin(nx)\bigr)$,则
$$\dfrac{1}{\pi}\int_{-\pi}^{\pi}f(x)^2\,\mathrm{d}x=\dfrac{a_0^2}{2}+\sum_{n=1}^{\infty}\bigl(a_n^2+b_n^2\bigr)$$.
\end{ti}

\textbf{证明}:记$S_n$为$f$的第$n$个$Fourier$部分和.由三角函数的正交性,对任意次数不超过$n$的三角多项式$T$,有
$$\int_{-\pi}^{\pi}(f-T)^2\,\mathrm{d}x=\int_{-\pi}^{\pi}(f-S_n)^2\,\mathrm{d}x+\int_{-\pi}^{\pi}(S_n-T)^2\,\mathrm{d}x.$$
因此$S_n$是此类多项式中平方积分误差最小的一个.下面说明任给$\varepsilon>0$,存在三角多项式$T$,使$\int_{-\pi}^{\pi}(f-T)^2<\varepsilon$.设$|f|\leq M$.由$f$的$Riemann$可积性,可取分割$P$,使$\sum_i\operatorname{osc}_{I_i}(f)|I_i|$任意小.在每个小区间$I_i$上取$f$的一个值构成阶梯函数$s$,则
$$\int_{-\pi}^{\pi}(f-s)^2\,\mathrm{d}x\leq2M\sum_i\operatorname{osc}_{I_i}(f)|I_i|.$$
只在阶梯函数有限个跳跃点附近的总长度很小的区间内将其改成线性连接,并在$-\pi,\pi$附近同样连接两端,可得连续的周期分段线性函数$g$,使$\int(s-g)^2$任意小.由定理8.7,可取三角多项式$T$使$\lVert g-T\rVert_\infty$任意小.利用$(f-T)^2\leq3[(f-s)^2+(s-g)^2+(g-T)^2]$,即得所需逼近.于是对充分大的$n$有$\int(f-S_n)^2<\varepsilon$,从而$\int(f-S_n)^2\to0$.最后由正交性,
$$\frac{1}{\pi}\int_{-\pi}^{\pi}f^2\,\mathrm{d}x-\left[\frac{a_0^2}{2}+\sum_{k=1}^n(a_k^2+b_k^2)\right]=\frac{1}{\pi}\int_{-\pi}^{\pi}(f-S_n)^2\,\mathrm{d}x\longrightarrow0,$$
即得$Parseval$恒等式.

\begin{ti}
\textbf{推论8.12(广义$Parseval$恒等式)}  设$f,g\in R[-\pi,\pi]$,其$Fourier$系数分别为$\{a_n,b_n\}$,$\{\alpha_n,\beta_n\}$,则
$$\frac{1}{\pi}\int_{-\pi}^{\pi}f(x)g(x)\,\mathrm{d}x=\frac{a_0\alpha_0}{2}+\sum_{n=1}^{\infty}(a_n\alpha_n+b_n\beta_n).$$
\end{ti}

\textbf{证明}:分别对$f+g$和$f-g$应用$Parseval$恒等式,然后二者相减即可.

\begin{ti}
\textbf{定理8.13($Fourier$系数的唯一性)}  设函数$f,g$在$[-\pi,\pi]$上连续且具有相同的$Fourier$系数,则$f\equiv g$.
\end{ti}

\textbf{证明}:设$h=f-g$,则$h\in C^0[-\pi,\pi]$且其$Fourier$系数全为零.由$Parseval$恒等式,$\dfrac{1}{\pi}\displaystyle\int_{-\pi}^{\pi}h(x)^2\,\mathrm{d}x=0$.因$h^2$为连续非负函数,故$h\equiv 0$,即$f\equiv g$.
\BookEnd
